% Source: https://artofproblemsolving.com/community/c3h3055143_challenging_geo_problem
% Stable identity: aops-topic-3055143-q01
Use $ A=(0,0),B=(1,0),C=(1,1),D=(0,1)$, $ E=(1-t,0),F=(1,1-t)$. Write $ u=1-t $. For $ u>0 $, the lines are $ y=ux $ and $ y=1-\dfrac{x}{u}$, giving $ P=(\dfrac{u}{1+u^2},\dfrac{u^2}{1+u^2})$. The formula also holds at $ u=0 $, where $ P=A $. With base $ CD=1 $, the area is \[\boxed{[CPD]=\dfrac{1}{2(1+(1-t)^2)}}.\]As $ t $ varies from 0 to 1, the area ranges from $\boxed{\dfrac{1}{4}\text{ to }\dfrac{1}{2}}$. It is not determined by the source conditions alone.
\begin{asy}
usepackage("amsmath");
size(7cm);
pair A=(0,0),B=(1,0),C=(1,1),D=(0,1),Ee=(0.6,0),F=(1,0.6),P=(0.6/1.36,0.36/1.36);fill(C--P--D--cycle,lightgray);draw(A--B--C--D--cycle^^A--F^^D--Ee);label("$A$",A,SW);label("$B$",B,SE);label("$C$",C,NE);label("$D$",D,NW);label("$E$",Ee,S);label("$F$",F,E);label("$P$",P,S);
\end{asy}
